\editionPage[\pdfbookmark[0]{Fiches outils}{bm-outils-debut-parent}]{outils-debut}{Fiche outil}{Poser et contrôler les opérations}{À consulter au fil des chapitres.}
\begin{blocEdition}{outils-debut}
\editionEtape{Aligner}Unités sous unités, virgules alignées pour les décimaux, rangs vides complétés par des zéros.\editionEtape{Diviser}Encadrer par des multiples du diviseur, abaisser chiffre à chiffre, écrire 0 au quotient quand le chiffre abaissé ne suffit pas. Pour un quotient décimal, transformer le reste en dixièmes, puis en centièmes, puis en millièmes.\editionEtape{Parenthèses}Calculer d’abord l’intérieur des parenthèses, en commençant par les plus intérieures. Une ligne par étape.\editionEtape{Contrôler}Estimer avant ; vérifier après : soustraction par addition, division par multiplication.\editionRetenir{Une opération posée se termine par son contrôle. Un résultat loin de l’estimation se refait.}
\end{blocEdition}
\editionPage{outils-2}{Fiche outil}{Résoudre un problème en quatre étapes}{À consulter au fil des chapitres.}
\begin{blocEdition}{outils-2}
\editionEtape{Comprendre}Raconter l’histoire avec ses mots : ce que l’on sait, ce que l’on cherche. Les mots « plus » ou « moins » ne décident pas de l’opération.\editionEtape{Modéliser}Schéma en barres pour les parties et le tout ; deux barres alignées pour comparer ; « fois plus, fois moins » pour la proportionnalité.\editionEtape{Calculer}Calcul mental quand c’est raisonnable, calcul posé sinon ; garder les étapes visibles.\editionEtape{Répondre et vérifier}Phrase réponse avec l’unité ; vérifier la vraisemblance : un trajet de 3 000 km à pied en une journée signale une erreur.\editionRetenir{Comprendre, modéliser, calculer, répondre : quand on bloque, repérer l’étape qui coince.}
\end{blocEdition}
\editionPage{outils-3}{Fiche outil}{Écrire un programme de construction}{À consulter au fil des chapitres.}
\begin{blocEdition}{outils-3}
\editionEtape{Croquis codé}Dessiner la figure à main levée, nommer tous les points, coder les longueurs égales et les angles droits.\editionEtape{Choisir l’ordre}Commencer par l’objet dont tous les autres dépendent, souvent un segment de longueur connue. Ne jamais utiliser un point avant de l’avoir construit.\editionEtape{Écrire précisément}Une instruction par ligne, avec les mesures : « Trace le segment [AB] de 5 cm. », « Trace le cercle de centre A passant par B. »\editionEtape{Tester}Faire exécuter le programme par quelqu’un qui ne voit pas la figure ; comparer les deux tracés.\editionRetenir{Un bon programme s’exécute sans retour en arrière et ne laisse aucun choix au hasard.}
\end{blocEdition}
\editionPage{outils-4}{Fiche outil}{Lire un tableau ou un graphique}{À consulter au fil des chapitres.}
\begin{blocEdition}{outils-4}
\editionEtape{Observer d’abord}Titre, grandeurs, unités, total. Dans un tableau à double entrée, une case se lit au croisement d’une ligne et d’une colonne.\editionEtape{Lire les axes}Repérer le pas des graduations et l’origine de l’axe : un axe qui ne commence pas à zéro exagère les écarts.\editionEtape{Choisir le graphique}Courbe pour une évolution, barres pour comparer des effectifs, diagramme circulaire pour les parts d’un tout, tableau pour les valeurs exactes.\editionEtape{Croiser}Vérifier les totaux, confronter deux documents, toujours revenir aux nombres.\editionRetenir{Un graphique répond vite ; le tableau garde les valeurs exactes. Les deux se contrôlent l’un l’autre.}
\end{blocEdition}
